\newcommand{\haty}{\hat{\bm{y}}}
\newcommand{\corey}{\bm{y}}
% \newcommand{\state}{\bm{u}}
\newcommand{\nhost}{\hat{n}}
\newcommand{\ncore}{n}
\newcommand{\nstate}{m}

% Maps
\newcommand{\Ymap}{\mathcal{Y}}   % host -> core parameter map
\newcommand{\Mmap}{\mathcal{M}}   % core parameters -> state (FE solve)
\newcommand{\Tmap}{\mathcal{T}}   % composed map

\section{A Compositional Framework for Sensitivity Analysis}

% -------------------------------------------------------
\subsection{Motivation}
% -------------------------------------------------------
% Premise:
%   - DDM computes exact sensitivities of the discrete equilibrium,
%     but the quantities it differentiates with respect to are limited
%     to parameters that appear explicitly in the finite element
%     residual (material constants, nodal coordinates, applied loads, etc.).
%   - In practice, the design or uncertainty parameters of interest
%     are often higher-level quantities (e.g.\ cross-section geometry,
%     topology descriptors, stochastic field coefficients) that
%     determine the core model inputs through an arbitrarily complex
%     preprocessing chain.
%   - AD frameworks such as JAX can differentiate arbitrary
%     host-language code, but are unaware of the implicit function
%     defined by the finite element solve.
%   - Goal: compose DDM inside the FE kernel with AD outside it,
%     so that the user obtains end-to-end sensitivities with
%     respect to any host-level parameter.

% -------------------------------------------------------
\subsection{Setting}
% -------------------------------------------------------
Define three vector spaces and two maps.

  Host parameters:   $\haty \in \mathbb{R}^{\nhost}$
      — the quantities the analyst wishes to differentiate with
        respect to (section dimensions, load multipliers, random
        field coefficients, etc.).

  Core parameters:   $\corey \in \mathbb{R}^{\ncore}$
      — the quantities that appear explicitly in the finite element
        residual and to which DDM can be applied directly
        (Young's modulus, fiber coordinates and areas, etc.).

  Model state:        $\state \in \mathbb{R}^{\nstate}$
      — the discrete solution vector (e.g.\ nodal displacements)
        obtained by solving the governing equilibrium.

  Parameter map:
  \[
      \Ymap : \mathbb{R}^{\nhost} \to \mathbb{R}^{\ncore},
      \qquad \haty \mapsto \corey = \Ymap(\haty).
  \]
      — a user-defined, differentiable mapping implemented in host
        code (e.g.\ a Python/JAX function that converts section
        dimensions to fiber discretisations).  Because it is
        ordinary host code, its Jacobian
        $\partial\Ymap/\partial\haty$ is available through AD.

  Model map:
  \[
      \Mmap : \mathbb{R}^{\ncore} \to \mathbb{R}^{\nstate},
      \qquad \corey \mapsto \state = \Mmap(\corey).
  \]
      — the (implicit) map defined by the finite element analysis:
        given core parameters $\corey$, assemble and solve
        $\bm{R}(\state;\corey) = \bm{0}$ to obtain $\state$.
        Its Jacobian $\partial\Mmap/\partial\corey$ is furnished
        column-by-column by the DDM.

  Composed map:
  \[
      \Tmap = \Mmap \circ \Ymap :
      \mathbb{R}^{\nhost} \to \mathbb{R}^{\nstate},
      \qquad \haty \mapsto \Mmap\!\bigl(\Ymap(\haty)\bigr).
  \]

% -------------------------------------------------------
\subsection{Differential Operations}
% -------------------------------------------------------
Describe JVP and VJP of each map, then their composition.

\subsubsection{Jacobian--vector product (forward mode)}

  By the chain rule the JVP of the composed map is
  \[
      \mathrm{D}\Tmap(\haty)\,\dot{\haty}
      \;=\;
      \underbrace{%
          \mathrm{D}\Mmap\!\bigl(\Ymap(\haty)\bigr)
      }_{\text{kernel (DDM)}}
      \;\cdot\;
      \underbrace{%
          \mathrm{D}\Ymap(\haty)\,\dot{\haty}
      }_{\text{host (AD forward)}}.
  \]

  Concretely:
  \begin{enumerate}
      \item \textbf{Host forward pass.}
            Evaluate $\corey = \Ymap(\haty)$ and propagate the
            tangent $\dot{\corey} = \mathrm{D}\Ymap(\haty)\,\dot{\haty}$
            using AD in forward mode.
      \item \textbf{Kernel forward pass.}
            Solve the primal problem $\bm{R}(\state;\corey)=\bm{0}$
            and, for each core parameter $y_i$ with
            $\dot{y}_i \neq 0$, obtain the sensitivity
            $\partial\state/\partial y_i$ from the DDM.
      \item \textbf{Contraction.}
            Form the JVP:
            \[
                \dot{\state}
                = \sum_{i=1}^{\ncore}
                  \frac{\partial\state}{\partial y_i}\,\dot{y}_i.
            \]
  \end{enumerate}

\subsubsection{Vector--Jacobian product (reverse mode)}

  The VJP (adjoint) of the composed map is
  \[
      \bar{\haty}
      \;=\;
      \bigl[\mathrm{D}\Tmap(\haty)\bigr]^\top \bar{\state}
      \;=\;
      \underbrace{%
          \bigl[\mathrm{D}\Ymap(\haty)\bigr]^\top
      }_{\text{host (AD reverse)}}
      \;\cdot\;
      \underbrace{%
          \bigl[\mathrm{D}\Mmap(\corey)\bigr]^\top \bar{\state}
      }_{\text{kernel (DDM)}}.
  \]

  Concretely:
  \begin{enumerate}
      \item \textbf{Kernel adjoint pass.}
            Given the cotangent (objective gradient) $\bar{\state}$,
            compute
            $\bar{y}_i
              = \bigl(\partial\state/\partial y_i\bigr)^\top
                \bar{\state}$
            for each core parameter using the DDM sensitivities
            already available from the primal solve.
      \item \textbf{Host reverse pass.}
            Propagate $\bar{\corey}$ back through $\Ymap$ by
            invoking AD in reverse mode to obtain $\bar{\haty}$.
  \end{enumerate}

  Note: Because the DDM furnishes the full sensitivity matrix
  $\partial\state/\partial\corey$ column-by-column, both the JVP
  and VJP of $\Mmap$ are available at the same cost — the
  distinction between forward and reverse mode applies only to the
  host map $\Ymap$.

% -------------------------------------------------------
\subsection{Composition Protocol}
% -------------------------------------------------------
  The framework requires three ingredients, each supplied by a
  different layer:

  \begin{enumerate}
      \item \textbf{Kernel primitive.}
            The finite element library exposes $\Mmap$ as a
            primitive operation together with its JVP and VJP rules,
            both implemented via DDM.

      \item \textbf{Host parameter map.}
            The user writes $\Ymap$ in the host language.
            No special annotation is needed; the AD framework
            differentiates it automatically.

      \item \textbf{Automatic composition.}
            The AD framework composes the two layers by the chain
            rule.  From the user's perspective, the result is a
            single differentiable function
            $\Tmap : \haty \mapsto \state$ that can be passed to
            any gradient-based algorithm (optimisation, UQ,
            inverse problems) without further manual derivation.
  \end{enumerate}

% -------------------------------------------------------
\subsection{Illustrative Example}
% -------------------------------------------------------
  Consider a cantilever beam whose response $u$ at the tip is a
  function of:
    - beam length $L$,
    - Young's modulus $E$, and
    - flange width $b$ and web depth $d$ of a rectangular section.

  The core parameters are $\corey = (L,\, E,\, y_1, z_1, A_1,
  \ldots, y_N, z_N, A_N)$ — the coordinates and areas of $N$
  fibers that discretise the cross-section.

  The host parameters are $\haty = (L,\, E,\, b,\, d)$.

  The parameter map $\Ymap$ generates a fiber layout from the
  section dimensions:
  \[
      (b,d) \;\xmapsto{\;\Ymap\;}\;
      \{(y_k, z_k, A_k)\}_{k=1}^{N}.
  \]

  The model map $\Mmap$ solves the equilibrium and returns the tip
  displacement $u$.

  End-to-end gradients such as $\partial u/\partial d$ are then
  obtained automatically:
  \[
      \frac{\partial u}{\partial d}
      = \sum_{k=1}^{N}
        \biggl(
          \frac{\partial u}{\partial y_k}\frac{\partial y_k}{\partial d}
        + \frac{\partial u}{\partial z_k}\frac{\partial z_k}{\partial d}
        + \frac{\partial u}{\partial A_k}\frac{\partial A_k}{\partial d}
        \biggr),
  \]
  where the $\partial u/\partial(\cdot)$ terms come from DDM and
  the $\partial(\cdot)/\partial d$ terms come from AD of $\Ymap$.